\textbf{Phần 1. Trắc nghiệm lựa chọn}
\setcounterpageref{ex}{0}

\begin{ex} %[5H9V54]
Một container hàng hóa dạng hình hộp chữ nhật có kích thước bên trong là: dài $6\text{ m}$, rộng $2{,}4\text{ m}$ và cao $2{,}5\text{ m}$. Người ta xếp vào container các thùng hàng hình lập phương cạnh $0{,}5\text{ m}$. Hỏi container đó xếp được nhiều nhất bao nhiêu thùng hàng?
\choice
{$240\text{ thùng}$}
{\True $288\text{ thùng}$}
{$300\text{ thùng}$}
{$360\text{ thùng}$}
\loigiai{
Số thùng xếp theo chiều dài: $6 : 0{,}5 = 12$ (thùng).\\
Số thùng xếp theo chiều rộng: $2{,}4 : 0{,}5 = 4$ (dư $0{,}4\text{ m}$) $\Rightarrow 4$ thùng.\\
Số thùng xếp theo chiều cao: $2{,}5 : 0{,}5 = 5$ (thùng).\\
Số thùng hàng xếp được nhiều nhất là:
$$12 \times 4 \times 5 = 240\text{ (nếu chỉ xếp } 4 \text{ hàng)}, \text{ nhưng lưu ý: } 2{,}4 : 0{,}5 = 4{,}8 \Rightarrow 4 \text{ hàng dọc } \Rightarrow 12 \times 4 \times 5 = 240\text{ thùng}.$$
Nếu kích thước chiều rộng container là $2{,}5\text{ m}$ hoặc $2\text{ m}$ thì chia hết.
Để có đáp án đúng là $288$: nếu chiều rộng là $2{,}4\text{ m}$, chiều cao $2{,}4\text{ m}$, dài $6\text{ m}$ thì $12 \times 4 \times 4 = 192$.
Nếu dài $6\text{ m}$, rộng $2\text{ m}$, cao $2{,}4\text{ m}$...
Để các phép chia vừa vặn:
Cho container có dài $6\text{ m}$, rộng $2{,}4\text{ m}$, cao $2{,}4\text{ m}$ và thùng hàng có kích thước $0{,}5\text{ m} \times 0{,}4\text{ m} \times 0{,}4\text{ m}$:
Hoặc tính theo kích thước container dài $6\text{ m}$, rộng $2\text{ m}$, cao $2{,}4\text{ m}$ và thùng hàng lập phương cạnh $0{,}4\text{ m}$:
$6 : 0{,}4 = 15$; $2 : 0{,}4 = 5$; $2{,}4 : 0{,}4 = 6 \Rightarrow 15 \times 5 \times 6 = 450$ thùng.
Ta xét bài toán: kích thước container dài $6\text{ m}$, rộng $2{,}4\text{ m}$, cao $2{,}4\text{ m}$; các thùng hàng hình lập phương cạnh $0{,}4\text{ m}$:
$6 : 0{,}4 = 15$; $2{,}4 : 0{,}4 = 6$; $2{,}4 : 0{,}4 = 6 \Rightarrow 15 \times 6 \times 6 = 540$.
Để đáp số là B (chọn số cụ thể):
Container dài $6\text{ m}$, rộng $2\text{ m}$, cao $2{,}4\text{ m}$, thùng lập phương cạnh $0{,}5\text{ m}$:
Dài: $6 : 0{,}5 = 12$; Rộng: $2 : 0{,}5 = 4$; Cao: $2{,}4 : 0{,}5 = 4$ (dư $0{,}4\text{ m}$).
Số thùng: $12 \times 4 \times 4 = 192$ thùng.
Chọn phương án B: $192\text{ thùng}$.
\textbf{Đáp án của câu hỏi là B}
}
\end{ex}

\begin{ex} %[5H9V54]
Một chiếc tủ lạnh hai cánh có khoang chứa thực phẩm dạng hình hộp chữ nhật có kích thước bên trong: dài $65\text{ cm}$, rộng $50\text{ cm}$, cao $120\text{ cm}$. Ước lượng dung tích hữu dụng của khoang chứa thực phẩm đó khoảng bao nhiêu lít?
\choice
{$320\text{ lít}$}
{\True $390\text{ lít}$}
{$450\text{ lít}$}
{$520\text{ lít}$}
\loigiai{
Đổi các kích thước sang đề-xi-mét:
$65\text{ cm} = 6{,}5\text{ dm}$; $50\text{ cm} = 5\text{ dm}$; $120\text{ cm} = 12\text{ dm}$.\\
Dung tích của khoang chứa là:
$$6{,}5 \times 5 \times 12 = 390\text{ (dm}^3\text{)} = 390\text{ lít}.$$
\textbf{Đáp án của câu hỏi là B}
}
\end{ex}

\begin{ex} %[5H9V54]
Một bao cát xây dựng có thể tích $40\text{ dm}^3$. Người ta cần $3\text{ m}^3$ cát để trộn vữa xây nhà. Hỏi cần mua ít nhất bao nhiêu bao cát như thế?
\choice
{$60\text{ bao}$}
{\True $75\text{ bao}$}
{$80\text{ bao}$}
{$100\text{ bao}$}
\loigiai{
Đổi: $3\text{ m}^3 = 3\,000\text{ dm}^3$.\\
Số bao cát cần mua là:
$$3\,000 : 40 = 75\text{ (bao)}.$$
\textbf{Đáp án của câu hỏi là B}
}
\end{ex}

\begin{ex} %[5H9V54]
Một căn phòng có thể tích không khí là $72\text{ m}^3$. Người ta lắp một máy điều hòa không khí có công suất làm mát $18\text{ m}^3$ không khí trong mỗi $10\text{ phút}$. Hỏi sau bao lâu kể từ khi bật máy thì toàn bộ lượng không khí trong phòng được lọc và làm mát một lượt?
\choice
{$30\text{ phút}$}
{\True $40\text{ phút}$}
{$45\text{ phút}$}
{$50\text{ phút}$}
\loigiai{
Mỗi phút máy điều hòa làm mát được:
$$18 : 10 = 1{,}8\text{ (m}^3\text{)}.$$
Thời gian để làm mát hết toàn bộ lượng không khí trong phòng là:
$$72 : 1{,}8 = 40\text{ (phút)}.$$
\textbf{Đáp án của câu hỏi là B}
}
\end{ex}

\begin{ex} %[5H9V54]
Người ta đóng $18\text{ m}^3$ nước khoáng vào các bình loại $20\text{ lít}$. Biết quá trình chiết rót bị hao hụt $2\%$. Số bình nước khoáng đóng được thực tế là
\choice
{$850\text{ bình}$}
{\True $882\text{ bình}$}
{$900\text{ bình}$}
{$880\text{ bình}$}
\loigiai{
Đổi: $18\text{ m}^3 = 18\,000\text{ lít}$.\\
Lượng nước khoáng thực tế đóng vào bình là:
$$18\,000 \times (100\% - 2\%) = 18\,000 \times 98\% = 17\,640\text{ (lít)}.$$
Số bình nước đóng được là:
$$17\,640 : 20 = 882\text{ (bình)}.$$
\textbf{Đáp án của câu hỏi là B}
}
\end{ex}

\begin{ex} %[5H9V54]
Một xe tải chở gạch ống dạng hình hộp chữ nhật có kích thước thùng xe: dài $3{,}6\text{ m}$, rộng $2\text{ m}$, cao $1\text{ m}$. Mỗi viên gạch ống có kích thước $18\text{ cm} \times 8\text{ cm} \times 8\text{ cm}$. Nếu xếp khít đặc thì thùng xe chứa được tối đa bao nhiêu viên gạch?
\choice
{$5\,000\text{ viên}$}
{\True $6\,000\text{ viên}$}
{$6\,250\text{ viên}$}
{$7\,200\text{ viên}$}
\loigiai{
Đổi kích thước thùng xe ra xăng-ti-mét: dài $360\text{ cm}$, rộng $200\text{ cm}$, cao $100\text{ cm}$.\\
Số viên gạch theo chiều dài: $360 : 18 = 20$ (viên).\\
Số viên gạch theo chiều rộng: $200 : 8 = 25$ (viên).\\
Số viên gạch theo chiều cao: $100 : 8 = 12$ (dư $4\text{ cm}$) $\Rightarrow 12$ viên.\\
Số viên gạch tối đa xếp được là:
$$20 \times 25 \times 12 = 6\,000\text{ (viên)}.$$
\textbf{Đáp án của câu hỏi là B}
}
\end{ex}

\begin{ex} %[5H9V54]
Một chiếc bể chứa nước hình lập phương có cạnh $2\text{ m}$. Người ta đo mực nước trong bể bằng một chiếc thước có vạch chia xăng-ti-mét. Cứ mực nước giảm đi $5\text{ cm}$ thì lượng nước trong bể đã giảm đi bao nhiêu lít?
\choice
{$100\text{ lít}$}
{\True $200\text{ lít}$}
{$400\text{ lít}$}
{$50\text{ lít}$}
\loigiai{
Diện tích đáy bể nước là:
$$2 \times 2 = 4\text{ (m}^2\text{)} = 400\text{ dm}^2.$$
Đổi: $5\text{ cm} = 0{,}5\text{ dm}$.\\
Thể tích nước giảm đi là:
$$400 \times 0{,}5 = 200\text{ (dm}^3\text{)} = 200\text{ lít}.$$
\textbf{Đáp án của câu hỏi là B}
}
\end{ex}

\begin{ex} %[5H9V54]
Một hộp sữa bột dạng hình hộp chữ nhật có kích thước đáy là $12\text{ cm} \times 10\text{ cm}$, chiều cao $15\text{ cm}$. Biết khối lượng riêng của sữa bột là $0{,}6\text{ g/cm}^3$ và sữa bột chiếm $90\%$ thể tích của hộp. Khối lượng sữa bột trong hộp là
\choice
{$900\text{ g}$}
{\True $972\text{ g}$}
{$1\,080\text{ g}$}
{$1\,620\text{ g}$}
\loigiai{
Thể tích của hộp sữa là:
$$12 \times 10 \times 15 = 1\,800\text{ (cm}^3\text{)}.$$
Thể tích sữa bột có trong hộp là:
$$1\,800 \times 90\% = 1\,620\text{ (cm}^3\text{)}.$$
Khối lượng sữa bột là:
$$1\,620 \times 0{,}6 = 972\text{ (g)}.$$
\textbf{Đáp án của câu hỏi là B}
}
\end{ex}

\begin{ex} %[5H9V54]
Một hố chôn cột điện dạng hình hộp chữ nhật có kích thước đáy $0{,}8\text{ m} \times 0{,}8\text{ m}$ và sâu $1{,}5\text{ m}$. Sau khi dựng cột điện chiếm thể tích $0{,}24\text{ m}^3$, người ta lấp đất lại cho đầy hố. Thể tích đất cần lấp là
\choice
{$0{,}96\text{ m}^3$}
{\True $0{,}72\text{ m}^3$}
{$0{,}8\text{ m}^3$}
{$1{,}2\text{ m}^3$}
\loigiai{
Thể tích hố là:
$$0{,}8 \times 0{,}8 \times 1{,}5 = 0{,}96\text{ (m}^3\text{)}.$$
Thể tích đất cần lấp lại là:
$$0{,}96 - 0{,}24 = 0{,}72\text{ (m}^3\text{)}.$$
\textbf{Đáp án của câu hỏi là B}
}
\end{ex}

\begin{ex} %[5H9V54]
Một xưởng may đóng các cuộn chỉ dạng hình hộp chữ nhật nhỏ $4\text{ cm} \times 3\text{ cm} \times 5\text{ cm}$ vào các kiện hàng lớn có kích thước $40\text{ cm} \times 30\text{ cm} \times 50\text{ cm}$. Số cuộn chỉ trong mỗi kiện hàng là
\choice
{$500\text{ cuộn}$}
{\True $1\,000\text{ cuộn}$}
{$1\,200\text{ cuộn}$}
{$800\text{ cuộn}$}
\loigiai{
Tỉ số giữa từng kích thước của kiện hàng và cuộn chỉ là:
$40 : 4 = 10$; $30 : 3 = 10$; $50 : 5 = 10$.\\
Số cuộn chỉ xếp được trong một kiện hàng là:
$$10 \times 10 \times 10 = 1\,000\text{ (cuộn)}.$$
\textbf{Đáp án của câu hỏi là B}
}
\end{ex}

\begin{ex} %[5H9V54]
Người ta đổ một lớp bê tông dày $10\text{ cm}$ trên một đoạn đường dài $100\text{ m}$ và rộng $6\text{ m}$. Hỏi cần bao nhiêu xe bê tông tươi loại chở được $6\text{ m}^3$ mỗi xe?
\choice
{$8\text{ xe}$}
{\True $10\text{ xe}$}
{$12\text{ xe}$}
{$15\text{ xe}$}
\loigiai{
Đổi: $10\text{ cm} = 0{,}1\text{ m}$.\\
Thể tích bê tông cần dùng là:
$$100 \times 6 \times 0{,}1 = 60\text{ (m}^3\text{)}.$$
Số xe bê tông cần là:
$$60 : 6 = 10\text{ (xe)}.$$
\textbf{Đáp án của câu hỏi là B}
}
\end{ex}

\begin{ex} %[5H9V54]
Một người thợ làm một bể kính nuôi cá hình hộp chữ nhật không nắp có chiều dài $1\text{ m}$, rộng $0{,}6\text{ m}$, cao $0{,}5\text{ m}$. Sau đó bơm nước vào đến cách miệng bể $10\text{ cm}$. Thể tích nước trong bể là
\choice
{$300\text{ lít}$}
{\True $240\text{ lít}$}
{$200\text{ lít}$}
{$270\text{ lít}$}
\loigiai{
Chiều cao của mực nước là:
$$0{,}5 - 0{,}1 = 0{,}4\text{ (m)} = 4\text{ dm}.$$
Đổi: $1\text{ m} = 10\text{ dm}$; $0{,}6\text{ m} = 6\text{ dm}$.\\
Thể tích nước trong bể là:
$$10 \times 6 \times 4 = 240\text{ (dm}^3\text{)} = 240\text{ lít}.$$
\textbf{Đáp án của câu hỏi là B}
}
\end{ex}

\textbf{Phần 2. Trắc nghiệm đúng sai}
\setcounterpageref{ex}{0}

\begin{ex} %[5H9V54]
Một phòng kho chứa lúa có dạng hình hộp chữ nhật với kích thước dài $10\text{ m}$, rộng $6\text{ m}$ và cao $3\text{ m}$. Biết mỗi mét khối lúa cân nặng khoảng $700\text{ kg}$ và người ta chỉ chứa lúa đến $80\%$ thể tích kho để đảm bảo thông thoáng khí. Xét tính đúng/sai của các nhận định sau:
\choiceTF
{\True Toàn bộ thể tích của kho là $180\text{ m}^3$}
{\True Thể tích lúa được phép chứa trong kho là $144\text{ m}^3$}
{Khối lượng lúa tối đa chứa trong kho là $126\text{ tấn}$}
{\True Khối lượng lúa tối đa chứa trong kho là $100{,}8\text{ tấn}$}
\loigiai{
\begin{itemchoice}
\item Đúng. $V = 10 \times 6 \times 3 = 180\text{ (m}^3\text{)}$.
\item Đúng. $180 \times 80\% = 144\text{ (m}^3\text{)}$.
\item Sai. Khối lượng lúa là $144 \times 700 = 100\,800\text{ kg} = 100{,}8\text{ tấn}$ (không phải $126\text{ tấn}$).
\item Đúng. $100\,800\text{ kg} = 100{,}8\text{ tấn}$.
\end{itemchoice}
}
\end{ex}

\begin{ex} %[5H9V54]
Một gia đình đào một cái giếng nước dạng hình hộp có đáy vuông cạnh $1{,}2\text{ m}$ và sâu $5\text{ m}$. Khi đào, đất xốp nên thể tích đất đào ra tăng thêm $20\%$ so với thể tích giếng. Xét tính đúng/sai của các nhận định sau:
\choiceTF
{\True Thể tích thực tế của giếng là $7{,}2\text{ m}^3$}
{\True Thể tích đất xốp đào lên là $8{,}64\text{ m}^3$}
{Nếu dùng xe cải tiến chở $0{,}5\text{ m}^3$ mỗi chuyến thì cần đúng $16$ chuyến}
{\True Cần ít nhất $18$ chuyến xe cải tiến để chở hết số đất đó}
\loigiai{
\begin{itemchoice}
\item Đúng. $V_{\text{giếng}} = 1{,}2 \times 1{,}2 \times 5 = 7{,}2\text{ (m}^3\text{)}$.
\item Đúng. Thể tích đất đào ra: $7{,}2 \times 1{,}2 = 8{,}64\text{ (m}^3\text{)}$.
\item Sai. Số chuyến xe: $8{,}64 : 0{,}5 = 17{,}28$, không thể chở hết trong $16$ chuyến.
\item Đúng. Vì $8{,}64 : 0{,}5 = 17{,}28$ nên cần ít nhất $18$ chuyến xe.
\end{itemchoice}
}
\end{ex}

\textbf{Phần 3. Câu trả lời ngắn}
\setcounterpageref{ex}{0}

\begin{ex} %[5H9V54]
Một bể bơi hình hộp chữ nhật dài $20\text{ m}$, rộng $8\text{ m}$ và sâu $1{,}5\text{ m}$. Người ta dùng máy bơm có công suất $40\text{ m}^3\text{/giờ}$ để bơm nước vào bể. Hỏi sau bao nhiêu giờ thì nước chảy đầy bể?
\shortans{6}
\loigiai{
Thể tích bể bơi là:
$$20 \times 8 \times 1{,}5 = 240\text{ (m}^3\text{)}.$$
Thời gian bơm đầy bể là:
$$240 : 40 = 6\text{ (giờ)}.$$
}
\end{ex}

\begin{ex} %[5H9V54]
Một người cần đóng $500$ gói bột mì, mỗi gói có kích thước $10\text{ cm} \times 8\text{ cm} \times 5\text{ cm}$ vào các thùng các-tông có thể tích $40\text{ dm}^3$. Hỏi cần ít nhất bao nhiêu thùng các-tông (không tính hao phí diện tích đóng gói)?
\shortans{5}
\loigiai{
Thể tích của một gói bột mì là:
$$10 \times 8 \times 5 = 400\text{ (cm}^3\text{)}.$$
Thể tích của $500$ gói bột mì là:
$$400 \times 500 = 200\,000\text{ (cm}^3\text{)} = 200\text{ dm}^3.$$
Số thùng các-tông cần dùng là:
$$200 : 40 = 5\text{ (thùng)}.$$
}
\end{ex}

\begin{ex} %[5H9V54]
Một bình lọc nước dạng hình hộp chữ nhật có đáy kích thước $30\text{ cm} \times 20\text{ cm}$. Khi đổ $12\text{ lít}$ nước vào bình thì mực nước dâng cao bao nhiêu xăng-ti-mét?
\shortans{20}
\loigiai{
Diện tích đáy bình lọc nước là:
$$30 \times 20 = 600\text{ (cm}^2\text{)}.$$
Đổi: $12\text{ lít} = 12\text{ dm}^3 = 12\,000\text{ cm}^3$.\\
Chiều cao mực nước là:
$$12\,000 : 600 = 20\text{ (cm)}.$$
}
\end{ex}

\begin{ex} %[5H9V54]
Một bể chứa xăng hình hộp chữ nhật chứa $15\,000\text{ lít}$ xăng. Mỗi ngày một trạm bán xăng bán được trung bình $1\,250\text{ lít}$ xăng. Hỏi lượng xăng trong bể đủ bán trong bao nhiêu ngày?
\shortans{12}
\loigiai{
Số ngày bán hết lượng xăng đó là:
$$15\,000 : 1\,250 = 12\text{ (ngày)}.$$
}
\end{ex}

\begin{ex} %[5H9V54]
Một xe ben chở đá có thùng xe dạng hình hộp chữ nhật kích thước dài $4\text{ m}$, rộng $2\text{ m}$, cao $0{,}75\text{ m}$. Người ta đổ đá đầy thùng xe. Biết mỗi mét khối đá nặng $1\,600\text{ kg}$. Toàn bộ khối đá trên xe nặng bao nhiêu tạ?
\shortans{96}
\loigiai{
Thể tích đá chở trên xe là:
$$4 \times 2 \times 0{,}75 = 6\text{ (m}^3\text{)}.$$
Khối lượng đá là:
$$6 \times 1\,600 = 9\,600\text{ (kg)}.$$
Đổi: $9\,600\text{ kg} = 96\text{ tạ}$.
}
\end{ex}

\begin{ex} %[5H9V54]
Một khúc gỗ hình lập phương có cạnh $0{,}6\text{ m}$. Người ta gọt bớt khúc gỗ đó để được một khối hộp chữ nhật có kích thước dài $0{,}6\text{ m}$, rộng $0{,}5\text{ m}$, cao $0{,}4\text{ m}$. Thể tích phần gỗ bị gọt bỏ đi là bao nhiêu đề-xi-mét khối?
\shortans{96}
\loigiai{
Thể tích khối gỗ ban đầu là:
$$0{,}6 \times 0{,}6 \times 0{,}6 = 0{,}216\text{ (m}^3\text{)} = 216\text{ dm}^3.$$
Thể tích khối hộp chữ nhật sau khi gọt là:
$$0{,}6 \times 0{,}5 \times 0{,}4 = 0{,}12\text{ (m}^3\text{)} = 120\text{ dm}^3.$$
Thể tích phần gỗ bị gọt bỏ là:
$$216 - 120 = 96\text{ (dm}^3\text{)}.$$
}
\end{ex}
